\documentclass[10pt,a4paper]{amsart}
\usepackage[margin=19mm]{geometry}
\usepackage{amsmath,amssymb,mathtools}
\usepackage[scr=rsfs,cal=euler]{mathalfa}
\usepackage{xfrac}
\usepackage[unicode,hidelinks,bookmarks=false]{hyperref}
\newcommand{\ie}{i.e.\ }
\newcommand{\cf}{cf.\ }
\newcommand{\resp}{respectively\ }
\newcommand{\Subsection}{\subsection}
\providecommand{\nobreakdash}{\nobreak\hbox{-}}
\setlength{\emergencystretch}{2em}
\multlinegap0pt
\usepackage{bm}
\usepackage{bbm}
\usepackage{dsfont}
\usepackage{xspace}
\usepackage{cancel}
\usepackage{mathtools}
\usepackage{amsthm}
\makeatletter
\@ifundefined{theorem}{\newtheorem{theorem}{Theorem}}{}
\@ifundefined{lemma}{\newtheorem{lemma}[theorem]{Lemma}}{}
\makeatother
\makeatletter
\expandafter\ifx\csname proof\endcsname\relax
\newenvironment{proof}[1][Proof]{\textbf{#1.} }{\hfill\rule{0.5em}{0.5em}}
\fi
\makeatother
\title[Applications of optimal error bounds for some generalized tw]{Applications of optimal error bounds for some generalized two-step iterative processes in Banach spaces}
\author{Tan-Phuc Nguyen, Thai-Hung Nguyen, Tien-Khai Nguyen, Cong-Duy-Nguyen Nguyen, Trung-Hieu Huynh}
\date{Source: arXiv:2510.16403v1 (CC BY 4.0)}
\begin{document}
\maketitle

\noindent\emph{Source and licence.} Applications of optimal error bounds for some generalized two-step iterative processes in Banach spaces. Version arXiv:2510.16403v1 (CC BY 4.0), distributed under the Creative Commons Attribution 4.0 International licence, as recorded in the arXiv licence metadata for this exact version and shown on the abstract page https://arxiv.org/abs/2510.16403v1. Full rights notice: licenses/arxiv-2510.16403-CC-BY-4.0.txt.\par\smallskip

\noindent\textit{Editorial scope.} Counted unit: the Theorem whose source label is \texttt{theo-L-IG} in the article (source0.tex lines 225--239, source label \texttt{theo-L-IG}), delivered together with the article's own complete original proof of it (source0.tex lines 240--280, one \texttt{proof} environment of 41 lines). No mathematics is written, repaired or re-derived: statement and proof are the article's own text line for line. Required context retained uncounted: xn-IG (lines 141--143); ln-ineq (lines 154--156); lem2 (lines 159--166). Every definition, display and label of those lines is retained unchanged. Omitted from this package and not redistributed: 1--140, 144--153, 157--158, 167--224, 281--691. Nothing inside a retained excerpt is omitted or altered: every retained body is delivered exactly as the article's lines are written, so no omission receipt of any kind accompanies this package. Citations: the retained excerpts carry no citation command at all, so no bibliography entry of the article is transcribed and no reference is redistributed. Reference adaptation: none was needed and none was made; every reference inside the delivered unit resolves inside it, so no unresolved reference or citation is produced. Printed slips kept literal: no printed word, symbol, hypothesis or number of the counted statement or of its proof was altered. Layout and class: the article's own class is \texttt{article} with packages including amsmath, amssymb, esint, amscd, xspace, fancyhdr, color, authblk, srcltx, fontenc, bbm, amsxtra, latexsym, framed; the wrapper is the lane's pinned \texttt{amsart} editorial wrapper at 10pt on A4, which restates the theorem-like environments the retained text needs and adds the article's own macro definitions that the retained text needs (none), with extra wrapper packages (bm, bbm, dsfont, xspace, cancel, mathtools). The delivered numbering is the unit's own. Assets: no retained span contains a figure environment, a graphics-inclusion command, a PDF-page inclusion, a table or a TikZ picture; no image file is copied into this package, the package declares no asset dependency and no image file is redistributed with it. Licence: version arXiv:2510.16403v1 of this article is distributed under the Creative Commons Attribution 4.0 International licence, as recorded in the arXiv licence metadata for this exact version and shown on the abstract page https://arxiv.org/abs/2510.16403v1. A full rights notice is delivered at licenses/arxiv-2510.16403-CC-BY-4.0.txt. Characters: every retained excerpt is pure ASCII, so the delivered engine, XeLaTeX with its default Latin Modern text font, reports no missing character.
\begin{align}\tag{IG}\label{xn-IG}
\begin{cases} y_{n} = (1-a_n)x_n + a_nT_1(x_n), \\ x_{n+1} = (1-b_n)T_1(y_n) + b_nT_2(y_n); \end{cases}
\end{align}

\begin{align}\label{ln-ineq}
\dfrac{x}{x+1} \le \ln(1+ x) \le x.
\end{align}

\begin{lemma}\label{lem2}
Assume that sequence $(a_n) \subset [0, 1]$ and bounded sequence $(u_n)$  satisfies 
$$1 - u_na_n > 0, \quad \mbox{ for all } n\in \mathbb{N}.$$ 
Then, two series $\displaystyle \sum_{n \in \mathbb{N}} \dfrac{a_n}{1- u_na_n}$ and $\displaystyle \sum_{n \in \mathbb{N}} a_n$ jointly converge or jointly diverge. In this case, we will denote by
\begin{align*}
\sum_{n \in \mathbb{N}} \dfrac{a_n}{1- u_na_n} \sim \sum_{n \in \mathbb{N}} a_n.
\end{align*}
\end{lemma}

\begin{theorem}\label{theo-L-IG}
Assume that  
\begin{align}\label{need-IG}
			a_k < \dfrac{1}{1 + \kappa_1} \ \text{ and } \ b_k < \dfrac{\kappa_1}{\kappa_1 + \alpha_2}, \ \text{ for all } k \in \mathbb{N}.
		\end{align}			
Then the optimal lower bound of~\eqref{xn-IG} is determined by
			\begin{align}\label{Ln-IG}
				\mathcal{L}_n^{IG} := \displaystyle \prod_{k=0}^n  [1 - (1+ \kappa_1)a_k][\kappa_1 - (\kappa_1 + \alpha_2)b_k],
			\end{align}
			for all $n \in \mathbb{N}$. Moreover, 
			\begin{enumerate}
				\item [(i)] If $\kappa_1 < 1$ then $\displaystyle \lim\limits_{n\to \infty} \mathcal{L}_n^{IG} = 0$ for every $(a_n), (b_n)$ satisfy~\eqref{need-IG}.
				\item [(ii)] If $\kappa_1 = 1$ then we have $\displaystyle \lim\limits_{n\to \infty} \mathcal{L}_n^{IG} = 0 \Longleftrightarrow \sum_{n \in \mathbb{N}} (a_n + b_n) = +\infty$. 
			\end{enumerate}
		\end{theorem}

		\begin{proof}
			Similarly, since $T_1 \in \text{Ne}_{\kappa_1, \alpha_1}(\Omega)$, $T_2 \in \text{Ne}_{\kappa_2, \alpha_2}(\Omega)$ and the parameter sequences $(a_n)$, $(b_n)$ satisfy the condition~\eqref{need-IG}, we can also have the following estimates
			\begin{align*}
				\| y_n - x^* \| &\ge (1 - a_n)\|x_n - x^*| - a_n\|T_1(x_n) - T_1(x^*)\| \\
				&\ge [1 - (1 + \kappa_1)a_n]\|x_n - x^*|, ~ n \in \mathbb{N},
			\end{align*}
			and
			\begin{align*} 
				\|x_{n+1}-x^*\| &\ge (1 - b_n)\|T_1(y_n) - T_1(x^*)\| - b_n\|T_2(y_n) - T_2(x^*)\| \\ &\ge [(1 - b_n)\kappa_1-b_n\alpha_2]\|y_n-x^*\|, \\
				&\ge [\kappa_1 - (\kappa_1 + \alpha_2)b_n][1 - (1 + \kappa_1)a_n]\|x_n - x^*|, ~ n \in \mathbb{N}.
			\end{align*}
			It leads to $\|x_{n+1}-x^*\| \ge \mathcal{L}_n^{IG}\|x_1-x^*\|$, for all $n \in \mathbb{N}$,
			where $\mathcal{L}_n^{IG}$ is defined by $\eqref{Ln-IG}$. 	On the other hand, we can easily check that if we choose two mappings $T_1,T_2$ defined by
			\begin{align*}
				T_1(x)= \kappa_1 x, \ T_2(x)=-\alpha_2 x, \quad x \in B\left(0, \dfrac{1}{2}\right),
			\end{align*}
			then we can observe that $x^*=0$ and 
			$$\|x_{n+1}-x^*\|=\mathcal{L}_n^{IG}\|x_0-x^*\|,$$
			for every $n \in \mathbb{N}$. Therefore, $\mathcal{L}_n^{IG}$ is the optimal lower bound of the iteration~\eqref{xn-IG}. \\
			
			In the first case $\kappa_1<1$, for all $k \in \mathbb{N}$, we have
			\begin{align*}
				[\kappa_1 - (\kappa_1 + \alpha_2)b_k][1 - (1 + \kappa_1)a_k] \le \kappa_1,~ \text{for all} ~k \in \mathbb{N}.
			\end{align*}
			It leads to $\mathcal{L}_n^{IG}\le \kappa_1^n$. Thus, we deduce that $\lim{\mathcal{L}_n^{IG}}=0$. \\
			
			For the second case $\kappa_1=1$, for every $n \in \mathbb N^*$, $\ln \mathcal{L}_n^{IG}$ can be rewritten as below
			\begin{align*}
				\ln{\mathcal{L}_n^{IG}} = \sum_{k=0}^{n} \ln[1 - (1 + \alpha_2)b_n] + \sum_{k=0}^{n} \ln(1 - 2a_n), ~ n \in \mathbb{N}.
			\end{align*}
By the inequality~\eqref{ln-ineq}, we have two following estimates
			\begin{align*}
				&\ln{\mathcal{L}_n^{IG}} \le -(1+\alpha_2)\sum_{k=0}^{n}{b_k}-2\sum_{k=0}^{n}{a_k},\\
				&\ln{\mathcal{L}_n^{IG}} \ge -(1+\alpha_2)\sum_{k=0}^{n} \dfrac{b_k}{1-(1+\alpha_2)b_k} - 2\sum_{k=0}^{n} \dfrac{a_k}{1-2a_k}.
			\end{align*}
Moreover, Lemma~\ref{lem2} gives us
			\begin{align*} 
				\sum_{k \in \mathbb{N}} \dfrac{b_k}{1-(1+\alpha_2)b_k} + \sum_{k \in \mathbb{N}} \dfrac{a_k}{1-2a_k}  \sim \sum_{k \in \mathbb{N}}a_k+ \sum_{k \in \mathbb{N}}b_k.
			\end{align*}
It allows us to conclude the second statement of this theorem. 
		\end{proof}

\end{document}
